\begin{proposition}[I6: closed-image exactness of residual columns; Proved]\label{i:closedexact}
Let $R$ be a regular noetherian local ring of dimension $d$,
$\mathfrak p$ a nonmaximal prime of height $h$, and $K_R$ its
residual resolution. With finite topologies and differentials by
postcomposition, for $j>h$ one has
\[
 \overline{\operatorname{im}\bigl(
  \Hom(E_{\mathfrak p},K_R^{j-1})\to
  \Hom(E_{\mathfrak p},K_R^j)\bigr)}
 =\ker\bigl(\Hom(E_{\mathfrak p},K_R^j)\to
             \Hom(E_{\mathfrak p},K_R^{j+1})\bigr).
\]
The kernel at $j=h$ is zero.
\end{proposition}
\begin{proof}
Continuity gives the inclusion from left to right. For the other
inclusion, take a cycle $f:E_{\mathfrak p}\to K_R^j$ and put
finitely many prescribed inputs in $Re$ with $e=e_n/u$ as in
Proposition~\ref{i:adjacent}. Its annihilator $I$ is
$\mathfrak p$-primary and
$I R_{\mathfrak p}=(x_1^n,\ldots,x_h^n)$.
The finite module $M=\Ext_R^j(R/I,R)$ vanishes after localization
at $\mathfrak p$ when $j>h$, by the length-$h$ Koszul resolution
there. Thus some $v\notin\mathfrak p$ annihilates $M$.

The element $e/v\in E_{\mathfrak p}$ has the same annihilator
$I$. The restriction of $f$ to $R(e/v)\simeq R/I$ is a cycle.
Its restriction to $Re$ multiplies the cohomology class by $v$,
because $f(e)=v f(e/v)$. This class is zero. Therefore some map
$b:Re\to K_R^{j-1}$ has differential equal to $f$ on $Re$.
Its cyclic source implies that $b$ has finitely many nonzero
components. Extend each component to $E_{\mathfrak p}$ by
injectivity and sum, obtaining
$g:E_{\mathfrak p}\to K_R^{j-1}$. Its differential agrees with
$f$ on the prescribed inputs, proving density in the kernel.

At $j=h$, only the diagonal summand receives a map from
$E_{\mathfrak p}$, so the Hom group is $\widehat{R_{\mathfrak p}}$.
Choose a prime immediately above $\mathfrak p$. The injectivity
in Proposition~\ref{i:adjacent} gives zero kernel for the next
differential.
\end{proof}

For a height-gap-two interval $\mathfrak p\subset\mathfrak q$,
apply this with $R=A_{\mathfrak q}$ and $j=h+1$. It identifies
the kernel of postcomposition from the full adjacent-arrow family
to $\Hom(E_{\mathfrak p},E_{\mathfrak q})$ as the closure of
the family obtained from the preceding residual differential and
$\widehat{A_{\mathfrak p}}$. This is an exact relation statement
in the actual Hom topologies. It does not identify the closure
inside an independently completed free algebra. The quotient-topology
and universal-map continuity gate therefore remains.
